% ============================================================
\section{为轨分析准备设计与输入数据}
\subsection*{Preparing Design and Input Data for Rail Analysis}

在使用 RedHawk/RedHawk-SC Fusion 能力分析设计中的压降与电流密度违例之前，使用 \cmd{set\_app\_options} 指定库与所需输入文件的位置。

表~\ref{tab:rail-appopts} 列出用于在 Fusion Compiler 环境中指定设计与输入数据以运行 RedHawk/RedHawk-SC 轨分析的应用选项。

\begin{table}[htbp]
\centering
\caption{用于指定设计与输入数据的应用选项（原书 Table 66）}
\label{tab:rail-appopts}
\small
\begin{tabularx}{\textwidth}{@{}l X@{}}
\toprule
\textbf{应用选项} & \textbf{说明} \\
\midrule
\appopt{rail.lib\_files} &
指定 \texttt{.lib} 格式的库文件。例如：
\begin{lstlisting}
fc_shell> set_app_options \
     -name rail.lib_files \
     -value {test1.lib test2.lib}
\end{lstlisting} \\
\appopt{rail.tech\_file} &
指定用于 PG 抽取与 PG 电迁移分析的 Apache 技术文件。例如：
\begin{lstlisting}
fc_shell> set_app_options \
   -name rail.tech_file \
   -value ChipTop.tech
\end{lstlisting} \\
\appopt{rail.3dic\_enable\_die} &
指定纳入详细分析的 die。例如：
\begin{lstlisting}
fc_shell> set_app_options \
     -name rail.3dic_enable_die \
     -value {DIE1 true DIE2 true ...}
\end{lstlisting} \\
\appopt{rail.instance\_power\_file} &
指定用户定义的 Instance Power File（IPF），用于向 pin 提供功耗。例如：
\begin{lstlisting}
fc_shell> set_app_options \
     -name rail.instance_power_file \
     -value {E1 FILE1 DIE2 FILE2 ...}
\end{lstlisting} \\
\appopt{rail.apl\_files} &
指定 Apache APL 文件，包含用于动态轨分析的电流波形与本征寄生。例如：
\begin{lstlisting}
fc_shell> set_app_options \
   -name rail.apl_files -value \
   {cell.current current \
   cell.cdev cap}
\end{lstlisting} \\
\appopt{rail.switch\_model\_files} &
指定用于动态轨分析的 RedHawk 开关单元模型文件。例如：
\begin{lstlisting}
fc_shell> set_app_options \
   -name rail.switch_model_files \
   -value {switch.model}
\end{lstlisting} \\
\appopt{rail.macro\_models} &
指定宏模型文件。例如：
\begin{lstlisting}
fc_shell> set_app_options \
   -name rail.macro_models \
   -value {gds_cell1 mm_dir1 \
   gds_cell2 mm_dir2}
\end{lstlisting} \\
\appopt{rail.lef\_files} &
（可选）指定 LEF 文件列表。未指定时，工具使用设计库中的数据生成一份。 \\
\appopt{rail.def\_files} &
（可选）指定 DEF 文件列表。未指定时，工具使用设计库中的数据生成一份。 \\
\appopt{rail.pad\_files} &
（可选）指定 pad 位置文件列表。未指定时，需通过运行 \cmd{create\_taps} 指定理想电压源。 \\
\appopt{rail.sta\_file} &
（可选）指定静态时序文件（STA），包含最终 slew 与延迟信息。未指定时，工具使用设计库中的数据生成静态时序窗信息。 \\
\appopt{rail.spef\_files} &
（可选）指定 SPEF 文件列表，包含信号 net 的详细寄生阻性与容性负载数据。未指定时，Fusion Compiler 使用设计库中的数据生成 SPEF 文件。 \\
\appopt{rail.effective\_resistance\_instance\_file} &
（可选）指定实例文件，列出用于有效电阻计算的单元实例。 \\
\appopt{rail.generate\_file\_type} &
（可选）指定为 RedHawk Fusion 或 RedHawk-SC Fusion 轨分析生成的脚本文件类型。可用类型：
\begin{itemize}
  \item \texttt{tcl}（默认）：生成 Tcl 格式文件，默认文件名为 \texttt{set\_user\_input.tcl}。
  \item \texttt{python}：生成 Python 格式文件，默认文件名为 \texttt{input\_files.py}。Python 脚本文件仅在 RedHawk-SC Fusion 轨分析流程中支持。
\end{itemize}
使用定制 Python 脚本文件运行 RedHawk-SC 轨分析时需要此应用选项。详见「支持 RedHawk-SC 定制 Python 脚本文件」。例如：
\begin{lstlisting}
fc_shell> set_app_options \
  -name rail.generate_file_type -value python
\end{lstlisting} \\
\appopt{rail.generate\_file\_variables} &
仅在 RedHawk-SC 中支持。指定定制 Python 脚本用于 RedHawk-SC 轨分析时的首选用户变量。当脚本包含非标准 collateral Python 变量名时需要此选项。支持的文件类型：LEF、DEF、SPEF、TWF、PLOC。例如：
\begin{lstlisting}
fc_shell> set_app_options \
   -name rail.generate_file_variables -value \
   {PLOC ploc LEF lef_files DEF def_files \
   SPEF spef_files TWF tw_files}
\end{lstlisting} \\
\appopt{rail.enable\_new\_rail\_scenario} &
（可选）启用带多个轨场景的 RedHawk Fusion 或 RedHawk-SC Fusion 轨分析。默认为 \texttt{false}。详见「使用多个轨场景运行轨分析」。 \\
\appopt{rail.enable\_parallel\_run\_rail\_scenario} &
（可选）在分析流程内启用分布式 RedHawk Fusion 或 RedHawk-SC Fusion 分析。 \\
\appopt{rail.scenario\_name} &
（可选）指定用于 RedHawk 或 RedHawk-SC Fusion 轨分析的场景。工具在电源完整性流程的 IR 驱动功能中遵从所指定的设计场景，例如 IR 驱动布局、IR 驱动并发时钟与数据优化以及 IR 驱动优化。详见「使用多个轨场景运行轨分析」。 \\
\appopt{rail.dump\_icc2\_results\_style} &
指定如何以新的 Ansys 地图样式加载结果（例如 \texttt{-value 2.0}）。 \\
\bottomrule
\end{tabularx}
\end{table}

\subsection{生成轨分析脚本文件}
\subsubsection*{Generating Rail Analysis Script Files}

可运行 \cmd{analyze\_rail} \opt{-script\_only}，根据设计库中的信息生成运行轨分析所需的设置，而不实际执行分析。

该命令将数据写入工作目录。保存到目录的输出文件见表~\ref{tab:script-only-out}。

\begin{table}[htbp]
\centering
\caption{运行 \texttt{analyze\_rail -script\_only} 时的输出数据（原书 Table 67）}
\label{tab:script-only-out}
\small
\begin{tabularx}{\textwidth}{@{}X X c c@{}}
\toprule
\textbf{文件名} & \textbf{说明} & \textbf{RH} & \textbf{RH-SC} \\
\midrule
RedHawk GSR 文件 & 包含运行 RedHawk Fusion 轨分析的配置设置 & 是 & 否 \\
LEF/DEF、SPEF 与 STA 文件 & 若运行 \cmd{analyze\_rail} 前未指定这些输入文件则生成 & 是 & 是 \\
RedHawk Fusion 运行脚本（\texttt{analyze\_rail.tcl}） & 包含运行轨分析所需命令 & 是 & 否 \\
RedHawk-SC Fusion 运行脚本（\texttt{analyze\_rail.py}） & 包含运行轨分析所需的 Python 命令 & 否 & 是 \\
\bottomrule
\end{tabularx}
\end{table}

\subsection{支持 RedHawk-SC 定制 Python 脚本文件}
\subsubsection*{Supporting RedHawk-SC Customized Python Script Files}

RedHawk-SC Fusion 提供应用选项与命令以读入用户定制的 Python 脚本文件。工具从设计库生成全部设计 collateral（包括 DEF、LEF、SPEF、TWF 与 PLOC），并将其加入用户定制 Python 脚本，以便在单次 \cmd{analyze\_rail} 迭代中运行压降分析。用户定制 Python 脚本文件支持多轨场景的轨分析，以及其它 IR 驱动功能（例如 IR 驱动布局或 IR 驱动并发时钟与数据优化）。

要为 RedHawk-SC Fusion 轨分析使用定制 Python 脚本文件：

\begin{enumerate}
  \item 将脚本文件类型设为 \texttt{python}：
\begin{lstlisting}
fc_shell> set_app_options -name rail.generate_file_type -value python
\end{lstlisting}
  \item 设置 \appopt{rail.generate\_file\_variables}，指定要加入用户定制 Python 脚本的输入文件（包括 DEF、LEF、SPEF、TWF 与 PLOC）的变量。若定制 Python 脚本使用与 RedHawk-SC Fusion 相同的值，则无需设置此应用选项。

\begin{center}
\small
\begin{tabular}{@{}lllll@{}}
\toprule
 & DEF & LEF & SPEF & TWF / PLOC \\
\midrule
\texttt{fc\_shell} 默认名 & DEF & LEF & SPEF & TWF / PLOC \\
RedHawk-SC Fusion 默认变量名 & \texttt{def\_files} & \texttt{lef\_files} & \texttt{spef\_files} & \texttt{tw\_files} / \texttt{ploc} \\
\bottomrule
\end{tabular}
\end{center}

例如，若为 DEF 变量指定 \texttt{test\_def\_file} 以匹配定制 Python 脚本中的 DEF 文件：
\begin{lstlisting}
fc_shell> set_app_options -name rail.generate_file_variables \
   -value {DEF test_def_files}
\end{lstlisting}
  \item 创建并指定轨场景，以生成要加入定制 Python 脚本的必要输入文件。

运行 \cmd{create\_rail\_scenario} 创建轨场景并将其与当前设计中的设计场景关联。可创建多个轨场景并与同一设计场景关联。

然后运行 \cmd{set\_rail\_scenario} \opt{-generate\_file} \texttt{\{LEF | TWF | DEF | SPEF | PLOC\}}，指定轨场景以及要从设计库生成的输入文件类型。
\end{enumerate}

\begin{noteBox}
使用多角多模技术运行轨分析时，必须在 \cmd{set\_rail\_scenario} 之前使用 \cmd{create\_rail\_scenario} 创建轨场景。
\end{noteBox}

下例创建 \texttt{T\_low} 与 \texttt{T\_high} 轨场景并为其生成输入文件：
\begin{lstlisting}
fc_shell> create_rail_scenario -name T_low \
   -scenario func_cbest
fc_shell> set_rail_scenario -name T_low \
   -generate_file {PLOC DEF SPEF TWF LEF} \
   -custom_script_file ./customized.py
fc_shell> create_rail_scenario -name T_high \
   -scenario func_cbest
fc_shell> set_rail_scenario -name T_high \
   -generate_file {PLOC DEF SPEF TWF LEF} \
   -custom_script_file ./customized.py
\end{lstlisting}

随后可将 \texttt{customized.py} 用于多轨场景轨分析或其它 IR 驱动功能。

\subsubsection{示例}
\subsubsection*{Examples}

\textbf{网格系统上的多轨场景轨分析（定制 Python 脚本）：}
\begin{lstlisting}
#Required. Specify the below options to enable multi-rail-scenario rail analysis.
set_app_options -list {
            rail.enable_new_rail_scenario true
            rail.enable_parallel_run_rail_scenario true
}

#Required. Set the type of the script file to python.
set_app_options -name rail.generate_file_type -value python

#Use user-defined test_def_files for DEF variable to match
#with customized python.
set_app_options -name rail.generate_file_variables \
   -value {DEF test_def_files }

create_rail_scenario -name func.ss_125c -scenario func.ss_125c
set_rail_scenario -name func.ss_125c -voltage_drop dynamic

create_rail_scenario -name func.ff_125c_ -scenario func.ff_125c
set_rail_scenario -name func.ff_125c \
   -generate_file {PLOC DEF SPEF TWF LEF} \
   -custom_script_file ./customized.py

set_host_options -submit_protocol sge \
   -submit_command {qsub -V -notify -b y \
   -cwd -j y -P bnormal -l mfree=16G}

analyze_rail -rail_scenario {func.ss_125c func.ff_125c} \
   -submit_to_other_machine
\end{lstlisting}

\textbf{本地机器上的多轨场景轨分析：}
\begin{lstlisting}
set_app_options -list {
            rail.enable_new_rail_scenario true
            rail.enable_parallel_run_rail_scenario true
}
set_app_options -name rail.generate_file_type -value python
set_app_options -name rail.generate_file_variables \
   -value {DEF test_def_files }

create_rail_scenario -name func.ss_125c -scenario func.ss_125c
set_rail_scenario -name func.ss_125c -voltage_drop dynamic

create_rail_scenario -name func.ff_125c_ -scenario func.ff_125c
set_rail_scenario -name func.ff_125c \
   -generate_file {PLOC DEF SPEF TWF LEF} \
   -custom_script_file ./customized.py

set_host_options -submit_command {local}

analyze_rail -rail_scenario {func.ss_125c func.ff_125c}
\end{lstlisting}

\textbf{网格系统上的单场景轨分析：}
\begin{lstlisting}
set_app_options -list {
            rail.enable_new_rail_scenario true
            rail.enable_parallel_run_rail_scenario true
}
set_app_options -name rail.generate_file_type -value python
set_app_options -name rail.generate_file_variables \
   -value {DEF test_def_files }

create_rail_scenario -name func.ff_125c_ -scenario func.ff_125c
set_rail_scenario -name func.ff_125c \
   -generate_file {PLOC DEF SPEF TWF LEF} \
   -custom_script_file ./customized.py

set_host_options -submit_protocol sge \
   -submit_command {qsub -V -notify -b y \
   -cwd -j y -P bnormal -l mfree=16G}

analyze_rail -rail_scenario {func.ff_125c} \
   -submit_to_other_machine
\end{lstlisting}

\textbf{网格系统上的 IR 驱动布局：}
\begin{lstlisting}
set_app_options -list {
            rail.enable_new_rail_scenario true
            rail.enable_parallel_run_rail_scenario true
}
set_app_options -name rail.generate_file_type -value python
set_app_options -name rail.generate_file_variables \
   -value {DEF test_def_files }

create_rail_scenario -name func.ss_125c -scenario func.ss_125c
set_rail_scenario -name func.ss_125c -voltage_drop dynamic

create_rail_scenario -name func.ff_125c_ -scenario func.ff_125c
set_rail_scenario -name func.ff_125c \
   -generate_file {PLOC DEF SPEF TWF LEF} \
   -custom_script_file ./customized.py

set_host_options -submit_protocol sge \
   -submit_command {qsub -V -notify -b y \
   -cwd -j y -P bnormal -l mfree=16G}
# Enable IR-driven placement before running clock_opt -from final_opto
set_app_options -name opt.common.power_integrity -value true

clock_opt -from final_opto
\end{lstlisting}

\textbf{多角多模设计上网格系统中的并发时钟与数据优化：}
\begin{lstlisting}
#Required. Specify the below options to enable MCMM-based optimization.
set_app_options -list {
            rail.enable_new_rail_scenario true
            rail.enable_parallel_run_rail_scenario true
}
set_app_options -name rail.generate_file_type -value python
set_app_options -name rail.generate_file_variables \
   -value {DEF test_def_files }

create_rail_scenario -name func.ss_125c -scenario func.ss_125c
set_rail_scenario -name func.ss_125c -voltage_drop dynamic

create_rail_scenario -name func.ff_125c_ -scenario func.ff_125c
set_rail_scenario -name func.ff_125c \
   -generate_file {PLOC DEF SPEF TWF LEF} \
   -custom_script_file ./customized.py

set_host_options -submit_protocol sge \
   -submit_command {qsub -V -notify -b y -cwd -j y \
   -P bnormal -l mfree=16G}
# Enable IR-driven CCD analysis before running route_opt
set_app_options -name opt.common.power_integrity -value true
route_opt
\end{lstlisting}

\begin{seeAlsoBox}
指定 RedHawk 与 RedHawk-SC 工作目录
\end{seeAlsoBox}

% ============================================================
\section{将理想电压源指定为 Tap}
\subsection*{Specifying Ideal Voltage Sources as Taps}

Tap 用于对被分析器件所处的外部电压源环境建模；它们不是设计本身的一部分，可视为电压源的虚拟模型。设计中理想电压源与理想电源的位置是获得准确轨分析结果所必需的。

要创建 tap，使用 \cmd{create\_taps}。可通过多次运行 \cmd{create\_taps} 创建多个 tap。

使用 \opt{-name} 为所创建的 tap 指定名称。若指定名称已存在于设计中，工具发出警告并用新 tap 替换原 tap。当通过单次调用 \cmd{create\_taps} 创建多个 tap 时，所有创建的 tap 按 \texttt{tap\_name\_NNN} 命名约定命名，其中 \texttt{NNN} 为唯一整数标识符，\texttt{tap\_name} 为 \opt{-name} 的字符串参数。

可用下列任一方式创建 tap：

\begin{itemize}
  \item \textbf{将顶层 PG pin 当作 tap}\\
    当 block 没有可用的物理 pad 或 bump 信息时，使用 \opt{-top\_pg}。

    使用 \opt{-supply\_net} 将 tap 显式关联到从其定义对象所连接的电源（电源或地）net。要隐式关联，运行不带 \opt{-supply\_net} 的 \cmd{create\_taps}。

    下例为所有顶层电源 pin 创建 tap：
\begin{lstlisting}
fc_shell> create_taps -supply_net VDD -top_pg
\end{lstlisting}

  \item \textbf{在绝对坐标处创建 tap}\\
    使用 \opt{-layer} 与 \opt{-point} 在指定坐标创建 tap。若指定位置没有电源 net、电源 pin 或 via 形状（即没有到电源网络的导电通路），则尽管 tap 存在，也对轨分析无影响。

    若在已有 tap 的位置再定义 tap，工具发出警告并用新 tap 替换原 tap。必须使用 \opt{-supply\_net} 指定该 tap 所关联的电源 net。

\begin{noteBox}
当 \opt{-point} 与 \opt{-of\_objects} 联用时，位置相对于指定对象的原点；否则相对于当前顶层 block 的原点。
\end{noteBox}

    下例使用绝对坐标创建 tap。命令检查该坐标位置是否触及任何布局形状：
\begin{lstlisting}
fc_shell> create_taps -supply_net VDD \
   -layer 3 -point {100.0 200.0}
Warning: Tap point (100.0, 200.0) on layer 3 for net VDD does not
touch any polygon (RAIL-305)
\end{lstlisting}

  \item \textbf{使用现有实例作为 tap}\\
    要在指定对象上创建 tap，使用 \cmd{create\_taps} \opt{-of\_objects}。对象列表中的每一项可以是单元或库单元，或单元与库单元的集合。默认情况下，工具在对象的电源与地 pin 上创建 tap。若指定 \opt{-point}，工具相对于每个对象原点在指定物理位置创建 tap。

    在分析已物理例化 pad cell 的顶层设计时，使用现有实例创建 tap。

    下例为 \texttt{some\_pad\_cell} 库单元的每个实例创建 tap。工具在相对单元实例原点 \texttt{\{30.4 16.3\}} 处的 metal1 层为每个实例创建一个 tap：
\begin{lstlisting}
fc_shell> create_taps -of_objects [get_lib_cell */some_pad_cell]\
   -layer metal1 -point {30.4 16.3}
\end{lstlisting}

    下例在匹配 \texttt{*power\_pad\_cell*} 模式的所有单元实例的全部电源与地 pin 上创建 tap：
\begin{lstlisting}
fc_shell> create_taps -of_objects [get_cells -hier \
   *power_pad_cell*]
\end{lstlisting}

  \item \textbf{导入 tap 文件}\\
    若处理的 block 尚未定义 pin，或 pad cell 的位置/设计尚未最终确定，可在 ASCII 文件中定义 tap 位置、层号与电源 net，再用 \cmd{create\_taps} \opt{-import} 导入。

    支持的 tap 文件格式如下：
\begin{lstlisting}
net_name layer_number [X-coord Y-coord]
\end{lstlisting}
    以分号（\texttt{;}）或井号（\texttt{\#}）开头的行视为注释。$x$/$y$ 坐标以微米为单位。

    例如：
\begin{lstlisting}
fc_shell> exec cat tap_file
   VDD 3 500.000 500.000
   VSS 5 500.000 500.000
fc_shell> create_taps -import tap_file
\end{lstlisting}

  \item \textbf{带 package 创建 tap}\\
    简单 package RLC（电阻--电感--电容）模型为所有电源与地 pad 提供一组固定的电气特性。要在轨分析中使用简单 package RLC 模型，需通过 \appopt{rail.pad\_files} 指定定义该简单 RLC 模型的 pad 位置文件。

    关于简单 package RLC 模型的详细说明，见 \textit{RedHawk User Manual} 相关章节。

\begin{noteBox}
在使用简单 package RLC 模型创建 tap 之前，必须使用 \appopt{rail.allow\_redhawk\_license\_checkout} 启用 RedHawk 签核许可证密钥。
\end{noteBox}
\end{itemize}

\subsection{验证 Tap 并查找无效 Tap}
\subsubsection*{Validating the Taps and Finding Invalid Taps}

创建 tap 时，默认 \cmd{create\_taps} 会验证所创建的 tap 是否插入在有效位置。若 tap 未触及任何布局形状，工具发出警告。要在形状外放置 tap 而不发出警告，使用 \opt{-nocheck} 禁用检查。

\begin{lstlisting}
fc_shell> create_taps -supply_net VDD -layer 3 \
   -point {100.0 200.0} -nocheck
\end{lstlisting}

\subsubsection{Tap 属性}
\subsubsection*{Tap Attributes}

验证检查期间，工具为每个插入的 tap 对象标记 \appopt{is\_valid\_location} 属性以指示其有效性。启用验证检查时，有效 tap 为 \texttt{true}，无效 tap 为 \texttt{false}。使用 \opt{-nocheck} 时，所有 tap 的属性值为 \texttt{unknown}。

工具在后续压降与最小路径电阻分析中更新 tap 属性。属于 PG net 的每个 tap 会对照抽取形状检查。分析完成后，工具根据轨分析结果更新 tap 有效性状态。

下例先列出本命令运行中加入 block 的 tap 对象属性，再检查 \texttt{Tap\_1437} 是否通过验证：
\begin{lstlisting}
fc_shell> list_attributes -application -class tap
...
Attribute Name            Object     Type       Properties  Constraints
-------------------------------------------------------------------------
context                   tap        string     A           integrity,
                                                            rail,
                                                            session
full_name                 tap        string     A
is_valid_location         tap        boolean    A
layer_name                tap        string     A
layer_number              tap        int        A
net                       tap        string     A
object_class              tap        string     A
parasitics                tap        string     A
position                  tap        coord      A
static_current            tap        double     A
type                      tap        string     A
  absolute_coord, auto_pg, cell_coord, cell_pg, lib_cell_coord,
  lib_cell_pg, signal, terminal, top_pg
fc_shell> get_attribute [get_taps Tap_1437] is_valid_location
true
\end{lstlisting}

\subsubsection{查找无效 Tap}
\subsubsection*{Finding Invalid Taps}

要查找未通过有效性检查的 tap，使用 \cmd{get\_taps} 报告 \appopt{is\_valid\_location} 为 \texttt{false} 的 tap。例如：
\begin{lstlisting}
fc_shell> get_attribute [ get_taps Tap_1437 ] is_valid_location
false
\end{lstlisting}

\subsubsection{为无效 Tap 查找替代位置}
\subsubsection*{Finding Substitute Locations for Invalid Taps}

若 tap 的指定位置无效，使用 \cmd{create\_taps} 的 \opt{-snap\_distance} 选项为 tap 查找替代位置。命令在同一层上（由 \opt{-layer} 定义）于指定 snap 距离内水平与垂直搜索附近布局形状。tap 被重新插入到距指定位置最小距离内的有效形状角点。

\begin{noteBox}
仅可为由 \opt{-point} 与 \opt{-layer} 指定的无效 tap 搜索替代位置。例如：
\begin{lstlisting}
fc_shell> create_taps -point {50,50} -layer {M1} \
   -snap_distance 100
Snap tap point from original location (50.000, 50.000) to new
 location (90.000, 65.910)
\end{lstlisting}
\end{noteBox}

如图~\ref{fig:tap-snap} 所示，指定的 tap 位置未触及任何布局形状。红色矩形是距指定位置最近的形状，其右下角在所有其它矩形中距离最小。工具将该 tap 重新插入到红色矩形的右下角。

\begin{figure}[htbp]
\centering
\fbox{\parbox{0.85\textwidth}{\centering\vspace{1.5cm}
\textit{（原书 Figure 201：Finding a Substitute Location for an Invalid Tap）}\\[0.3em]
\small Original \{x,y\} $\rightarrow$ Snapped tap。请参阅原版 PDF 插图。
\vspace{1.5cm}}}
\caption{为无效 Tap 查找替代位置}
\label{fig:tap-snap}
\end{figure}

\subsection{移除 Tap}
\subsubsection*{Removing Taps}

要移除先前用 \cmd{create\_taps} 创建的 tap，使用 \cmd{remove\_taps}。要移除特定 tap，对 \cmd{get\_taps} 使用 \opt{-filter}。下例移除类型为 \texttt{top\_pg} 的全部电源 pin tap：
\begin{lstlisting}
fc_shell> remove_taps [get_taps -filter type==top_pg]
\end{lstlisting}

要移除 block 中的全部 tap：
\begin{lstlisting}
fc_shell> remove_taps
\end{lstlisting}
